\chapter{Osteopenia}
\index{Osteopenia}

\medskip
\noindent\textit{Educational information; seek professional advice for individual fracture risk.}

\section*{What it means}
Osteopenia means bone mineral density is lower than the reference range for healthy young adult bone but not in the osteoporosis range. It is a measurement, not a complete prediction of fracture risk, and many people have no symptoms.

\section*{Assessment}
A DXA scan commonly measures bone density at the hip and spine. Clinicians also consider age, previous fractures, falls, medicines such as long-term steroids, family history, menopause or hormone status, smoking, alcohol, nutrition, and conditions that affect bone. A low result may need repeat testing or investigation for an underlying cause.

\section*{Reducing fracture risk}
Discuss weight-bearing and resistance exercise, balance and fall prevention, adequate nutrition, calcium and vitamin D intake, smoking, alcohol, vision, and home safety. Supplements are not automatically needed and should be individualized. Medicines are considered when overall fracture risk is high, not solely because the label says osteopenia.

\section*{When to seek medical care}
Arrange review after a fracture from a minor injury, repeated falls, height loss, new back pain, or long-term use of medicines that weaken bone. Seek urgent help after a major injury, inability to stand, severe new back or hip pain, or neurological symptoms.

\section*{Sources}
\begin{itemize}
\item NIAMS bone-density information: \url{https://www.niams.nih.gov/health-topics/bone-mineral-density-tests-what-numbers-mean}
\end{itemize}
